\documentclass[11pt]{article}
\usepackage[letterpaper,margin=0.78in]{geometry}
\usepackage{amsmath,amssymb,amsthm,mathtools}
\usepackage{graphicx,microtype,enumitem,xcolor,hyperref}
\hypersetup{colorlinks=true,linkcolor=blue!45!black,urlcolor=blue!55!black,citecolor=blue!45!black}
\setlist{nosep,leftmargin=1.6em}
\newtheorem{theorem}{Theorem}
\newtheorem{lemma}{Lemma}
\newcommand{\D}{\mathbb D}
\newcommand{\T}{\mathbb T}
\newcommand{\K}{\mathcal K}
\newcommand{\N}{\mathcal N}
\newcommand{\M}{\mathcal M}
\newcommand{\A}{A^2(\D)}
\definecolor{answer}{RGB}{120,35,45}

\title{A Fredholm-index obstruction to essential unitary equivalence\\
\large A negative answer to Question 6.4 of arXiv:2407.18455}
\author{Automated mathematical discovery run \texttt{fa\_banach\_001}, lane 9}
\date{17 August 2026}

\begin{document}
\maketitle

\begin{center}
\fcolorbox{answer}{answer!6}{\parbox{0.91\linewidth}{
\textbf{Candidate full counterexample.}
The two extensions in Question 6.4 need not be essentially unitarily
equivalent.  In the specialization \(p(z,w)=z-w\), \(B_1=B_2=\mathrm{id}\),
the second extension is the Bergman Toeplitz extension and its coordinate
lift has Fredholm index \(-1\).  For the first extension, the corresponding
coordinate lift is invertible and has index \(0\).  Essential unitary
equivalence would preserve this index.
}}
\end{center}

\section*{Source question and transcription}
The source chooses a distinguished homogeneous polynomial
\(p=\prod_i(z-\alpha_iw)^{n_i}\), finite Blaschke products \(B_1,B_2\), and
an infinite Blaschke product \(B\) whose zeros accumulate at every point of
\(\T\).  It sets
\[
 B_3=B\circ B_1,\qquad B_4=\prod_i B\circ(\alpha_iB_2),\qquad
 \M=[B_3(z)B_4(w),p_B].
\]
The paper proves that \(\M^\perp\) and \([p_B]^\perp\) are essentially
normal and have the same essential spectrum, hence produce two extensions
of the same commutative \(C^*\)-algebra by the compact operators.  Question
6.4 asks whether these extensions are essentially unitarily equivalent.

\section*{Definitions}
Two quotient-module extensions over \(C(X)\) are \emph{essentially unitarily
equivalent} if a unitary between their Hilbert spaces conjugates the lifts of
each coordinate function modulo compact operators.  In particular, whenever
a coordinate lift minus \(\lambda\) is Fredholm, its index is invariant under
such an equivalence.

The Bergman space \(\A\) has orthonormal basis
\(\{\sqrt{n+1}\,\zeta^n:n\geq0\}\).  Multiplication by \(\zeta\) is the
weighted shift with weights \(\sqrt{(n+1)/(n+2)}\); it is injective with
one-dimensional cokernel and therefore has index \(-1\).

\section*{Idea / proof intuition}
Collapse the bidisc quotient by \(z-w\) onto the diagonal.  The second
quotient becomes the Bergman module.  Adding \(B(z)B(w)\) becomes quotienting
the Bergman module by the invariant subspace generated by \(B^2\).  If
\(B(0)\ne0\), a one-line Bezout identity
\(a(z)z+B(0)^{-2}B(z)^2=1\) makes the new coordinate compression invertible.
The resulting index jump obstructs essential unitary equivalence.

\clearpage
\section*{Original source image}
\begin{center}
\includegraphics[width=0.98\textwidth,height=0.82\textheight,keepaspectratio]{source_question_crop.png}
\end{center}
\noindent\footnotesize Source: Wang--Zhao--Zhu, arXiv:2407.18455, PDF
p.~16, equations (6.1)--(6.2) and Question 6.4.  This is a direct crop from
\texttt{source\_paper.pdf}.

\clearpage
\section*{The specialized pair of quotient modules}
For \(m\geq2\) and \(0\leq k<m\), put
\[
 \beta_{m,k}=r_m e^{2\pi i k/m},\qquad
 r_m=1-\frac{2^{-m}}{m^2}.
\]
Then
\[
 \sum_{m=2}^\infty\sum_{k=0}^{m-1}(1-|\beta_{m,k}|)
 =\sum_{m=2}^\infty\frac{2^{-m}}m<\infty.
\]
Hence these points are the zeros of an infinite Blaschke product \(B\).
Their boundary cluster set is all of \(\T\), while \(B(0)\ne0\): the angular
grids become dense, and convergence of the displayed sum makes the product
of the zero moduli at the origin strictly positive.

Take
\[
 p(z,w)=z-w,\qquad B_1(z)=z,\qquad B_2(w)=w.
\]
This is an admissible instance of the source construction, with
\[
 p_B=z-w,\qquad B_3=B_4=B,
\]
and quotient modules
\[
 \N_0=[z-w]^\perp,qquad
 \N_1=[z-w,B(z)B(w)]^\perp.
\]
By the source result, both are essentially normal and
\(\sigma_e(\N_0)=\sigma_e(\N_1)=\{(\zeta,\zeta):\zeta\in\T\}\), which we
identify with \(\T\).

\begin{lemma}[Diagonal Bergman model]
There is a unitary \(U:\N_0\to\A\) under which both compressed coordinate
multipliers become \(M_\zeta\), and the compression of
\(B(z)B(w)\) becomes \(M_{B^2}\).  Moreover,
\[
 U\N_1=K:=\A\ominus\overline{B^2\A}.
\]
\end{lemma}

\begin{proof}
The vectors
\[
 e_n=\frac1{\sqrt{n+1}}\sum_{j=0}^n z^jw^{n-j},\qquad n\geq0,
\]
form an orthonormal basis of \([z-w]^\perp\).  Diagonal restriction sends
\(e_n\) to \(\sqrt{n+1}\,\zeta^n\), so it defines the asserted unitary.
Projection of either \(ze_n\) or \(we_n\) onto \(\N_0\) is
\(\sqrt{(n+1)/(n+2)}e_{n+1}\).  The bounded analytic functional calculus
then sends \(B(z)B(w)\) to \(B(M_\zeta)^2=M_{B^2}\).  Finally, quotienting
by the closed submodule generated by this multiplier gives the stated
orthogonal complement.  This complement is infinite-dimensional, since the
Bergman reproducing kernels at the infinitely many zeros of \(B\) belong to
it and are linearly independent.
\end{proof}

\clearpage
\section*{Index obstruction}
Let \(T=P_KM_\zeta|_K\) be the coordinate compression on \(K\).

\begin{theorem}
For the above admissible source data, the extensions associated with
\(\N_1\) and \(\N_0\) are not essentially unitarily equivalent.
\end{theorem}

\begin{proof}
Set \(c=B(0)^{-2}\) and
\[
 a(z)=\frac{1-cB(z)^2}{z}.
\]
The apparent singularity at zero is removable.  The resulting function is
in \(H^\infty\): it is bounded on \(|z|\leq1/2\) by compactness, and on
\(|z|\geq1/2\) by the displayed quotient.  Thus
\[
 a(z)z+cB(z)^2=1. \tag{1}
\]
Because \(\overline{B^2\A}\) is invariant under every bounded analytic
multiplier, compression to \(K\) defines a unital homomorphism
\(H^\infty\to B(K)\).  The compressed multiplier by \(B^2\) is zero.
Applying this homomorphism to (1) gives
\[
 a(T)T=Ta(T)=I_K.
\]
Hence \(T\) is invertible and \(\operatorname{ind}T=0\).

On \(\N_0\), the corresponding coordinate lift is unitarily equivalent to
the Bergman shift \(M_\zeta\), which has
\(\operatorname{ind}M_\zeta=-1\).  If the two extensions were essentially
unitarily equivalent, their coordinate lifts would be unitarily conjugate
modulo compact operators.  Fredholm index is unchanged by unitary
conjugacy and compact perturbation, forcing \(0=-1\), a contradiction.
\end{proof}

\paragraph{Stronger observation.}
The explicit radial-grid choice of \(B\) is only used to meet the source's
dense boundary-zero hypothesis and ensure \(B(0)\ne0\).  More generally,
for any admissible \(B\) and any \(\lambda\in\D\setminus Z(B)\), the same
argument with
\(c=B(\lambda)^{-2}\) and
\(a(z)=(1-cB(z)^2)/(z-\lambda)\) shows that \(T-\lambda\) is invertible,
whereas \(M_\zeta-\lambda\) has index \(-1\).

\clearpage
\section*{Verification, novelty, and scope}
\begin{itemize}
\item The Blaschke condition and full-circle boundary cluster set are
checked directly from the radial grid.
\item The diagonal basis calculation verifies the Bergman weights and the
index \(-1\) without invoking an abstract model theorem.
\item The passage from \(B(z)B(w)\) to \(B^2\) uses the quotient-module
homomorphism property, so no unproved equality of raw multiplier ranges is
needed; closures are retained throughout.
\item The Bezout coefficient \(a\) is explicitly bounded, avoiding any
appeal to the corona theorem.
\item Both left and right inverse identities hold because the compressed
\(H^\infty\) calculus is commutative.
\end{itemize}

\paragraph{Novelty audit.}
Cheap run indexes were searched for arXiv:2407.18455, its exact final
question, the two displayed extension formulas, and Bergman/model-space
index obstructions.  Bounded primary-source web/arXiv searches through 17
August 2026 found the source paper and later talks describing its main
results, but no later paper or preprint resolving Question 6.4.  Novelty
confidence is moderate rather than definitive.

\paragraph{Limitations.}
This gives a complete negative answer to the universal reading of Question
6.4 by exhibiting an admissible instance in which equivalence fails.  It
does not classify when other choices of \(p,B_1,B_2,B\) yield equivalent
extensions, and it does not resolve Yang's separate Hilbert--Schmidt
conjecture or the source's conjecture about all infinitely generated
submodules.

\paragraph{Human review.}
Pending.  Priority checks are the identification of the source's extension
symbol \(C(\sigma_e)\) with \(C(\T)\) in this specialization, and the
standard implication from essential unitary equivalence of extensions to
compact-unitary equivalence of the coordinate lifts.

\section*{References}
\begin{enumerate}[label={[\arabic*]}]
\item P. Wang, C. Zhao, and Z. Zhu, \emph{Essential normality of quotient
modules vs. Hilbert-Schmidtness of submodules in \(H^2(\D^2)\)},
arXiv:2407.18455, especially Section 6 and Question 6.4.
\item L. Brown, R. Douglas, and P. Fillmore, \emph{Extensions of
\(C^*\)-algebras and K-homology}, Ann. of Math. 105 (1977), 265--324.
\item R. G. Douglas, V. I. Paulsen, \emph{Hilbert Modules over Function
Algebras}, Longman Research Notes 217, 1989.
\end{enumerate}

\end{document}
